% Generated draft: requires mathematical and source review.
\documentclass[11pt]{article}
\usepackage[margin=25mm]{geometry}
\usepackage{fontspec}
\setmainfont{Latin Modern Roman}
\usepackage{amsmath,amssymb,mathtools,unicode-math}
\setmathfont{Latin Modern Math}
\usepackage{enumitem,needspace,xurl,hyperref}
\setlength{\parindent}{0pt}
\setlength{\parskip}{5pt}
\setlength{\emergencystretch}{3em}
\newcommand{\N}{\mathbb{N}}
\newcommand{\Z}{\mathbb{Z}}
\newcommand{\Q}{\mathbb{Q}}
\newcommand{\R}{\mathbb{R}}
\newcommand{\C}{\mathbb{C}}
\newcommand{\F}{\mathbb{F}}
\newcommand{\A}{\mathbb{A}}
\newcommand{\PP}{\mathbb{P}}
\newcommand{\E}{\mathbb{E}}
\newcommand{\eps}{\varepsilon}
\newcommand{\dd}{\,\mathrm{d}}
\begin{document}
\section*{UnsolvedMath ID: 6700004}
\textbf{Source problem code:} catalogue entry
\subsection{Definitions and mathematical statement}
For a Riemannian manifold $(M,g)$, scalar curvature is the metric trace $\operatorname{Sc}_g=\operatorname{tr}_g\operatorname{Ric}_g$. A lower bound written $\operatorname{Sc}_g\geq\sigma$ means pointwise comparison with the specified constant or function $\sigma$.

The source asks what such lower bounds imply for the topology and geometry of maps between manifolds. In a representative notation one considers maps $f:(M,g)\to(N,h)$ and seeks restrictions relating curvature hypotheses to the map's topological data and metric behavior. The record does not specify whether the bound is on the source, target, or both, the allowed manifolds and maps, or any conclusion. Degree, homotopy, distance distortion, and area distortion are possible kinds of data to study, not source-stated requirements.

This is a broad comparison and obstruction programme rather than a single quantified assertion. Closedness/completeness, dimension, spin assumptions, target curvature or geometry, map regularity, nonzero-degree hypotheses, and contraction notions must be fixed in any precise special theorem. None is imposed on all maps by this short prompt. In particular, no universal rigidity conclusion follows merely from naming a scalar lower bound without those additional data.
\subsection{Short English statement}
Investigate how scalar-curvature lower bounds restrict maps between manifolds, both topologically and metrically. The source poses a programme without a specified map class or comparison conclusion.
\subsection{Sources}
\begin{itemize}
\item Local UnsolvedMath registry, permanent ID 6700004, exact source code catalogue entry, scalar-curvature item [?4]; proposed by Misha Gromov, 2017. Registry status partially solved, classification PARTIAL-PROGRESS.
\item Misha Gromov, \emph{101 Questions, Problems and Conjectures around Scalar Curvature} (2017), printed tag [?4], occurrence 4, PDF page 1, \url{https://www.ihes.fr/~gromov/wp-content/uploads/2018/08/101-problemsOct1-2017.pdf}. Supplied normalized pdf-text extraction accessed July 29, 2026; public author-hosted incomplete PDF. No PDF was accessed here.
\item The supplied summary invokes Llarull-type sphere comparison at scalar-curvature scale $n(n-1)$, Dirac/Plateau and families-index/$K$-area methods, spherical length comparisons, and macroscopic-dimension programmes. It names Llarull, ``Goette--Semmler'', and Gromov without precise publications or hypotheses. Its informal area-shrinking/``non-isometry'' description is not a reliable exact theorem statement and is not adopted as one.
\item The summary reports many partial rigidity/obstruction results and no general theory for arbitrary maps/targets. Provenance flags current status and statement-reuse terms NEEDS\_REVIEW; these are supplied assessments, not an independent literature check.
\end{itemize}
\end{document}
